% !PRINTED none
% Abschnitt-opening (part-title) page: no running-head folio visible on this
% scan. Carries the Abschnitt heading, §.175 heading, and body text.

\begin{center}
\gesperrt{Siebzehnter Abschnitt.}

\gesperrt{Algebraische Auflösung von Gleichungen.}

\rule{3cm}{0.4pt}

\S.\,175.

\gesperrt{Reduction der Gruppe durch reine Gleichungen.}
\end{center}

Eine der ältesten Fragen, an der sich vorzugsweise die
neuere Algebra entwickelt hat, ist die nach der sogenannten
\gesperrt{algebraischen Auflösung} der Gleichungen, worunter man eine
Darstellung der Wurzeln einer Gleichung durch eine Reihe von
Radicalen, oder die Berechnung durch eine endliche Kette von
Wurzelziehungen versteht. Auf diese Frage fällt von der Gruppentheorie das hellste Licht.

Präcisiren wir zunächst die Frage, um die es sich handelt,
so ist es offenbar die, ob und wie man den Körper $\Omega$ durch
successive Adjunction von Wurzelgrössen (Radicalen) so erweitern
kann, dass entweder alle oder wenigstens ein Theil der Wurzeln
im erweiterten Körper enthalten sind. Eine Wurzelgrösse ist
aber eine solche, die zwar nicht selbst, von der aber irgend eine
ganze Potenz in $\Omega$ enthalten ist, also, wenn $a$ eine Grösse in $\Omega$
ist, die Wurzel einer Gleichung von der Form

\[
y^m - a = 0,
\]

d. h. einer \gesperrt{reinen Gleichung}.

Soll eine \gesperrt{irreducible} Gleichung algebraisch auflösbar sein,
oder wenigstens eine oder einige algebraisch darstellbare Wurzeln
haben, so muss nach einer successiven Adjunction von Wurzeln
reiner Gleichungen in endlicher Anzahl die gegebene Gleichung
reducibel werden, da ein Theil ihrer Wurzeln in dem erweiterten
Körper enthalten sein soll. Da die anfängliche Gruppe $P$ transitiv
ist, so muss diese Gruppe schliesslich intransitiv werden, oder
%CONT>>
